% part: intuitionistic-logic
% chapter: propositions-as-types
% section: proof-terms

\documentclass[../../../include/open-logic-section]{subfiles}

\begin{document}

\olfileid{int}{pty}{ter}

\olsection{ترم‌های برهان}

فرض‌ها را در \Log{N1i} با متغیرهایی (برای نمونه، $!A^x$) به‌عنوان
برچسب ترخیص نشانه‌گذاری می‌کنیم. در \Log{N2i}، هر دنبالهٔ
$\Gamma \Sequent !A$ در بافت سمت چپ~$\Gamma$ فهرستی از جفت‌های
متغیر--!!{formula} دارد. بنابراین در \Log{N2i} هر دنباله در
!!a{proof} نه‌تنها اطلاعاتی دربارهٔ اینکه !!{proof} تا آن نقطه چه چیزی
را !!{prove} کرده است (یعنی~$!A$)، بلکه اطلاعات فرض‌های باز در هر
نقطه (یعنی~$\Gamma$) را نیز حمل می‌کند. اگر اطلاعات مربوط به اینکه
$!A$ \emph{چگونه} !!{prove} شده است نیز در هر دنباله درج شود چه؟ برای
این منظور می‌توان دستگاه \emph{برچسب‌خورده با ترم}~\Log{tN2i} را
تعریف کرد که دنباله‌هایش به شکل $\Gamma \Sequent N : !A$ هستند.
~‌$\Gamma$ و $!A$ مانند پیش‌اند: مجموعه‌ای از جفت‌های $x: !B$ و
!!{formula}یی که زیر-!!{proof} منتهی به دنباله نشان داده از فرض‌های
~$\Gamma$ نتیجه می‌شود. $N$ ترمی ویژه است که !!{proof} منتهی به
دنباله را رمزگذاری می‌کند.

ترم‌هایی که به هر دنباله اختصاص می‌دهیم از متغیرهای $x$، $y$، \dots،
و با نمادهای تابعی رمزگذار استنتاج‌های انجام‌شده تا آن نقطه ساخته
می‌شوند. برای هر قاعدهٔ استنتاج به نماد تابعی جداگانه‌ای نیاز داریم.
نمادهای تابعی متناظر با قواعد !!{introduction} «سازنده» و نمادهای
متناظر با قواعد !!{elimination} «ویرانگر» نامیده می‌شوند. هر نماد تابع
به تعداد مقدمات قاعده آرگومان می‌گیرد. برای نمونه، نمادهای تابعی
\Intro{\land} و \Elim{\land}[1] می‌توانند به‌ترتیب $c^\land$
(دوموضعی) و $d^\land_1$ (تک‌موضعی) باشند. این قواعد در \Log{tN2i}
چنین می‌شوند:
\[
\Axiom$\Gamma_1 \fCenter N: !A$
\Axiom$\Gamma_2 \fCenter M: !B$
\RightLabel{\Intro{\land}}
\BinaryInf$\Gamma_1, \Gamma_2 \fCenter c^\land(N, M) : !A \land !B $
\DisplayProof
\quad
\Axiom$\Gamma \fCenter N: !A \land !B$
\RightLabel{\Elim{\land}[1]}
\UnaryInf$\Gamma \fCenter d^\land_1(N) : !A$
\DisplayProof
\]
اگر قاعده‌ای استنتاجی فرضی را مرخص کند---یعنی مقدمه‌ای داشته باشد که
در آن !!{formula} برچسب‌دار~$x:!A$ آمده ولی این فرمول در بافت نتیجه
نباشد---باید این امر را در ترم برهان نیز نشان دهیم. نمادهای تابعی
وابسته به چنین قواعدی، متغیرهای متناظر را \emph{مقید} می‌کنند. برای
نمونه، قاعدهٔ \Intro{\lif} از مقدمهٔ
$x:!A, \Gamma \Sequent !B$ به نتیجهٔ
$\Gamma \Sequent !A \lif !B$ می‌رود. سازندهٔ $c^{\lif}$ یک آرگومان
می‌گیرد، اما باید نشان دهد که متغیر~$x$ مقید می‌شود. این راه بیانِ ترم
برهان برای مرخص‌شدن فرض دارای برچسب~$x$ است. برای نمونه، اگر ترم برهان
مقدمه~$N$ باشد، می‌توانیم ترم برهان نتیجه را $c^{\lif}([x]N)$ و قاعده
را چنین بنویسیم:
\[
\Axiom$x:!A, \Gamma \fCenter N: !B$
\RightLabel{\Intro{\lif}}
\UnaryInf$\Gamma \fCenter c^{\lif}([x]N) : !A \lif !B$
\DisplayProof
\]

برای آنکه بتوانیم !!a{proof} را به‌طور کامل بازسازی کنیم، به اطلاعات
بیشتری نیاز داریم. اگر !!{formula}s مقدمهٔ یک قاعده
!!{formula} نتیجه را کاملاً تعیین نکنند، باید آن اطلاعات را به ترم
بیفزاییم. این وضع در \Intro{\lif} رخ می‌دهد، پس
$c^{\lif}([x:!A]N)$ می‌نویسیم. در \Intro{\lor} و~\FalseInt نیز چنین
است؛ بنابراین از نمادهای تابعی‌ای استفاده می‌کنیم که آن اطلاعات را
حمل کنند، برای نمونه:
\[
\Axiom$\Gamma \fCenter N:!A$
\RightLabel{\Intro{\lor}[1]}
\UnaryInf$\Gamma \fCenter c^{\lor{!B}}_1(N):!A \lor !B$
\DisplayProof
\quad
\Axiom$\Gamma \fCenter N:\lfalse$
\RightLabel{\FalseInt}
\UnaryInf$\Gamma \fCenter d^\lfalse_{!A}(N):!A$
\DisplayProof
\]

با این ایده‌ها می‌توان دستگاه استنتاج طبیعی به سبک دنباله‌ای ساخت که
هر دنبالهٔ آن به شکل $\Gamma \Sequent N : !A$ باشد و در آن~$N$ یک
\emph{ترم برهان} است.

میان چنین ترم‌های برهانی و ترم‌های \emph{حساب لامبدای نوع‌دار} ارتباط
نزدیکی وجود دارد. برای رعایت این ارتباط، به جای نمادهای عمومی سازنده
و ویرانگر بالا از نمادهای رایج در ادبیات حساب لامبدا استفاده می‌کنیم.
برای نمونه، به جای $c^{\lif}([x:!A]N)$ می‌نویسیم
$\lambd[\typeof{x}{!A}][N]$، به جای $d^{\lif}(N,M)$ می‌نویسیم
$(NM)$، به جای $c^\land(N,M)$ می‌نویسیم $\pair{N}{M}$، و به جای
~$d^\land_i(N)$ می‌نویسیم $\proj{i}{N}$.

\begin{defn}
  ترم‌های برهان با قواعد زیر به‌طور استقرایی ساخته می‌شوند:
  \begin{enumerate}
  \item هر متغیر~$x$ یک ترم برهان است (اصول).
  \item اگر $x$ متغیر، $!A$ !!a{formula} و $N$ ترم برهان باشد،
    $\lambd[\typeof{x}{!A}][N]$ نیز ترم برهان است~(\Intro\lif).
  \item اگر $N$ و $M$ ترم برهان باشند، $(NM)$ نیز ترم برهان
    است~(\Elim\lif).
  \item اگر $N$ و $M$ ترم برهان باشند، $\pair{N}{M}$ ترم برهان
    است~(\Intro\land).
  \item اگر $N$ ترم برهان باشد، $\proj{i}{N}$ نیز ترم برهان است، که
    در آن $i$ برابر $1$ یا~$2$ است~(\Elim\land[i]).
  \item اگر $N$ ترم برهان و $!A$ یک فرمول باشد،
    $\inj[!A]{i}{N}$ ترم برهان است، که در آن $i$ برابر $1$ یا~$2$
    است~(\Intro\lor[i]).
  \item اگر $M$، $N_1$ و $N_2$ ترم برهان، و $x_1$ و $x_2$ متغیر
    باشند، $\dcase{M}{x_1}{N_1}{x_2}{N_2}$ ترم برهان
    است~(\Elim\lor).
  \item اگر $N$ ترم برهان و $!A$ !!a{formula} باشد،
    $\abort{!A}{N}$ یک ترم برهان است~(\FalseInt).
  \end{enumerate}
\end{defn}

هر ترم برهانی ترم برهان یک !!a{proof} نیست. برای نمونه، ترم
$\pair{N}{M}$ برهانی را رمزگذاری می‌کند که به استنتاج~\Intro{\land}
پایان می‌یابد، پس !!{formula} پایانی آن یک عطف~$!A \land !B$ است.
این فرمول نمی‌تواند مقدمهٔ اصلی قاعدهٔ~\Elim{\lif} باشد؛ بنابراین
$(\pair{N}{M}P)$ نمی‌تواند ترم برهان یک برهان درست باشد. تنها ترم‌های
برهانی که از قواعد \Log{tN2i} پیروی می‌کنند ترم‌های برهان
\emph{درست}اند. این قواعد در \olref{tab:tN2ip} آمده‌اند.

\subfile{rules-tN2}

\end{document}
